\documentclass[10pt,a4paper]{amsart}
\usepackage[margin=19mm]{geometry}
\usepackage{amsmath,amssymb,mathtools}
\usepackage[scr=rsfs,cal=euler]{mathalfa}
\usepackage{xfrac}
\usepackage[unicode,hidelinks,bookmarks=false]{hyperref}
\newcommand{\ie}{i.e.\ }
\newcommand{\cf}{cf.\ }
\newcommand{\resp}{respectively\ }
\newcommand{\Subsection}{\subsection}
\providecommand{\nobreakdash}{\nobreak\hbox{-}}
\setlength{\emergencystretch}{2em}
\multlinegap0pt
\usepackage{amssymb}
\usepackage{mathtools}
\newtheorem{lemma}{Lemma}
\newtheorem{proposition}{Proposition}
\title{Equidistribution in shrinking sets for arithmetic spherical harmonics}
\author{Maximiliano Sanchez Garza}
\date{Source: arXiv:2603.00790v2 (CC BY 4.0)}
\begin{document}
\maketitle

\noindent\emph{Source and licence.} Equidistribution in shrinking sets for arithmetic spherical harmonics. Version arXiv:2603.00790v2 (CC BY 4.0), distributed by arXiv under the Creative Commons Attribution 4.0 International licence, http://creativecommons.org/licenses/by/4.0/, as recorded in the arXiv licence metadata for this exact version and shown on the abstract page https://arxiv.org/abs/2603.00790v2. Full licence notice: \texttt{licenses/arxiv-2603.00790-CC-BY-4.0.txt}.\par\smallskip

\noindent\textit{Editorial scope.} Counted unit: the article's proposition VarianceFormula (source0.tex lines 722--724). The article's complete original proof of it is delivered with it (source0.tex lines 726--739, one \texttt{proof} environment of 14 lines). No mathematics is written, repaired or re-derived: the statement and the proof are the article's own lines, in the article's own notation, and no retained line is substituted or re-flowed.  Retained uncounted as the source-local context the proof needs: lines 178--180; lines 480--482.  Nothing was omitted from any retained span, so the package carries no omission record.  No retained line contains a citation, so the wrapper carries no bibliography and no bibliography entry.  No retained line contains a figure environment, an image-inclusion command or drawing code, so no image is delivered and the package declares no asset dependency.  Editorial formatting only, in addition: an AMS \texttt{amsart} preamble that restates the article's own declarations and macros used by the retained lines, namely the environments \texttt{lemma}, \texttt{proposition} on separate counters, together with the standard packages those lines need.  The retained environments print in this document as Lemma 1, Proposition 1 in order of appearance, whereas the article prints the same items with the numbering of its own full document.  The complete native source of the article as distributed in the arXiv e-print for this exact version is bundled unchanged as \texttt{sources/195513/source0.tex}.
\begin{lemma}\label{SpectralRes}
Let $$ f_0(z) := \frac{1}{\sqrt{\mathrm{vol}(\mathcal{O}^{\times} \backslash S^2)}},$$ so that $\langle f_0, f_0 \rangle = 1$, and let $\mathcal{B}$ be an orthonormal basis of spherical harmonics of $L^2(\mathcal{O}^{\times} \backslash S^2)$. Then a function $g \in L^2(\mathcal{O}^{\times} \backslash S^2)$ has the spectral expansion, valid in the $L^2$-sense, of the form $$ g(z) = \langle g, f_0 \rangle f_0(z) + \sum_{f\in \mathcal{B}} \langle g,f\rangle f(z).$$ Moreover, Parseval's identity holds: for $g_1, g_2 \in L^2(\mathcal{O}^{\times} \backslash S^2)$, we have $$ \langle g_1, g_2 \rangle = \langle g_1, f_0\rangle \langle f_0, g_2\rangle + \sum_{f\in \mathcal{B}} \langle g_1, f\rangle \langle f, g_2\rangle.$$
\end{lemma}

\begin{lemma}\label{IntBallEqualInnerProd}
If $\phi : \mathcal{O}^{\times} \backslash S^2 \rightarrow \mathbb{C}$ is a spherical harmonic of Laplacian eigenvalue $\lambda_{\psi}^2 = \ell_{\psi}(\ell_{\psi}+1)$, then $$ \frac{1}{\mathrm{vol}(B_R)} \int_{B_R(w)} \phi(z) \: \mathrm{d}\sigma(z) = \langle \phi, K_R(\cdot, w) \rangle = h_R(\ell_{\phi}) \phi(w).$$
\end{lemma}

\begin{proposition}\label{VarianceFormula}
Let $\psi$ be an $L^2$-normalized spherical harmonic of Laplacian eigenvalue $\lambda_{\psi}^2 = \ell_{\psi}(\ell_{\psi}+1)$ for some integer $\ell_{\psi}>0$. Then, for $R>0$, $$\mathrm{Var}(\psi;R) = \sum_{\phi \in \mathcal{B}} |h_R(\ell_{\phi})|^2 \left| \langle |\psi|^2, \phi \rangle \right|^2.$$
\end{proposition}

\begin{proof}
By Lemmas \ref{SpectralRes} and \ref{IntBallEqualInnerProd}, we have 
\begin{align*}
    \langle |\psi|^2, K_R(\cdot, w) \rangle &= \langle |\psi|^2, f_0\rangle \langle f_0, K_R(\cdot, w)\rangle + \sum_{\phi \in \mathcal{B}} \langle |\psi|^2, \phi\rangle \langle \phi, K_R(\cdot, w)\rangle \\
    &= \frac{1}{\mathrm{vol}(\mathcal{O}^{\times} \backslash S^2 )} + \sum_{\phi \in \mathcal{B}} h_R(\ell_{\phi}) \phi(w) \langle |\psi|^2 , \phi \rangle.
\end{align*}
This equation implies, again by Lemma \ref{IntBallEqualInnerProd}, that $$ \frac{1}{\mathrm{vol}(B_R)} \int_{B_R(w)} |\psi(z)|^2 \: \mathrm{d}\sigma(z) - \frac{1}{\mathrm{vol}(\mathcal{O}^{\times} \backslash S^2)} = \sum_{\phi \in \mathcal{B}} h_R(\ell_{\phi}) \phi(w) \langle |\psi|^2 , \phi \rangle.$$ Since $\mathcal{B}$ is an orthonormal basis of $L^2(\mathcal{O}^{\times} \backslash S^2)$, squaring and integrating over $w$ we have
\begin{align*}
    \mathrm{Var}(\psi;R) &= \int_{\mathcal{O}^{\times} \backslash S^2} \left| \sum_{\phi \in \mathcal{B}} h_R(\ell_{\phi}) \phi(w) \langle |\psi|^2 , \phi \rangle \right|^2 \mathrm{d}\sigma(w) \\
    &= \int_{\mathcal{O}^{\times} \backslash S^2} \sum_{\phi_1, \phi_2 \in \mathcal{B}} h_R(\ell_{\phi_1}) \phi_1(w) \langle |\psi|^2 , \phi_1 \rangle \overline{h_R(\ell_{\phi_2}) \phi_2(w) \langle |\psi|^2 , \phi_2 \rangle} \: \mathrm{d}\sigma(w) \\
    &= \sum_{\phi_1, \phi_2 \in \mathcal{B}} h_R(\ell_{\phi_1}) \overline{h_R(\ell_{\phi_2})}  \langle |\psi|^2 , \phi_1 \rangle \overline{\langle |\psi|^2 , \phi_2 \rangle} \int_{\mathcal{O}^{\times} \backslash S^2} \phi_1(w) \overline{\phi_2(w)} \: \mathrm{d}\sigma(w) \\
    &= \sum_{\phi \in \mathcal{B}} |h_R(\ell_{\phi})|^2 \left| \langle |\psi|^2, \phi \rangle \right|^2. \qedhere
\end{align*}
\end{proof}

\end{document}
